\section{Conclusion}

We proved that a common clustered multiplier enters an orthogonal two-lane
covariance through its squared modulus, not its outer sign.  A complete
nonorthogonal expansion turns every possible real-sign effect into a graph-cut
edge, and hard-window synthesis limits common-sign variation by
$2\eps\norm W\norm B$.  The result closes one tempting route: outer $C_h$
signs cannot control the same-bucket main term.  The open road is now the local
lane covariance and its literal V59 attachment, beginning with a sharp
longitudinal--transverse decomposition.
